\scanpage{279}
\mspage{229}
d'autre par\unsure{t}
\diagram{\begin{tikzcd}[ampersand replacement=\&]
X\times I \arrow[r,"\mathrm{pr}_1"] \& X
\end{tikzcd}}
est $W_{\unsure{\infty|w}}$-parfait, d'où la factorisation cherchée
\diagram{\begin{tikzcd}[ampersand replacement=\&]
X\amalg X \arrow[r,"i_X"] \& X\times I \arrow[r,"p_X=\mathrm{pr}_1^X"] \& X
\end{tikzcd}}
-- il n'y a pas plus simple ! Comment l'utiliser pour construire la factorisation d'une flèche $f:\unsure{\alpha|X}\to Y$ quelconque, en dualisant les constructions relatives aux espaces de chemins ? On définit $Z(f)$ par le carré cocartésien dans $\mathrm{Cat}$
\diagram{\begin{tikzcd}[ampersand replacement=\&,column sep=large]
X \arrow[r,dotted,"{\deleted{Z(f)}}","{\mathrm{inj}_{\unsure{1|2}}^X}"'] \& X\times I \arrow[r,"f'"] \& Z(f) \arrow[d,bend left=42,dashed,"{\unsure{P|p}_f}"]\\
\& X \arrow[u,hook,"{\mathrm{inj}_0^X}"'] \arrow[r,"f"] \& Y \arrow[u,hook,"{\unsure{w|W}}"']
\end{tikzcd},}
on a $Z(f)\xrightarrow{\alpha}Y$ induit par $\mathrm{pr}_1:X\times I\longrightarrow X$,
\margin{\figure{\begin{tikzpicture}[x=1ex,y=1ex]
\draw (0,2.9) -- (.16,2.75) -- (.20,2.55);
\draw (.56,2.67) .. controls (.95,2.69) and (.74,2.45) .. (.91,2.33) -- (.85,2.23)
.. controls (.98,1.87) and (.80,1.72) .. (.60,1.67) -- (.32,.66) -- (1.22,1.05);
\end{tikzpicture}}}
et la factorisation cherchée de $f$ est
\diagram{\begin{tikzcd}[ampersand replacement=\&,column sep=large]
X \arrow[r,"{i_f=f'\circ\mathrm{inj}_1^X}"] \& Z(f)\arrow[r,"P_f"] \& Y.
\end{tikzcd}}
Je dis que
\[
\left\{\begin{aligned}
i_f&\text{ est imm. ouverte de Dwyer, donc une}\\[-2pt]
&\qquad W\text{-cofibration}\\
p_f&\ \deleted{\unsure{\text{est }W\text{-asphérique}|\text{est }W\text{-parfait}}}\in W
\end{aligned}\right.
\]
% archive sheet number: 273
